%% Ported from drafts/methods_v05.md on 2026-09-21 (markdown to LaTeX; prose unchanged except em-dash removal).
%% This file is now the working copy: edit here, not in the markdown draft. Cite keys are from references.bib.

\section{Methods}\label{sec:methods}

\subsection{Design overview: the RUME and the Geography, Population, Movement decomposition}

The central design goal of \texttt{epymorph} is to let a researcher change one assumption of
a spatial modeling experiment while holding everything else fixed, so that the effect of
that assumption can be isolated rather than confounded with incidental differences in code.
To make this possible, \texttt{epymorph} organizes a complete experiment around a single
object, the Runnable Modeling Experiment (RUME). A RUME contains everything required to
execute a simulation and nothing more: a model of disease transmission, a model of host
movement, a geographic scope, a time frame, an initialization scheme, and the parameters,
supplied as constants or functions. Because a RUME is fully self-contained, it can be saved,
shared, and re-executed without reference to external scripts, which directly addresses the
reproducibility problems that motivate the package~\citep{oshinubi_systematic_2026}.

Internally, a RUME realizes the Geography, Population, Movement (GPM) framework introduced
above~\citep{oshinubi_systematic_2026}: the intra-population model (IPM), the movement model (MM), and the geographic
model (GEO), represented in software as a geographic scope together with the data attached
to its nodes, are separate, interchangeable components. In a single-strata experiment these
are combined directly. In a multi-strata experiment, the components that may differ between
strata (the IPM, the MM, and the initialization scheme) are grouped into one GPM object per
stratum, each representing a distinct subpopulation such as an age or risk group, while the
geographic scope and time frame are shared across strata; the strata are then composed into
a single model (see ``Multi-strata model composition'' below). The practical payoff is that
swapping an SIR process for an SEIR process, or replacing a gravity-style movement kernel
with empirical commuting flows, means changing a single component and leaving the rest of
the experiment untouched, turning comparisons that would otherwise require rewriting a
monolithic codebase into controlled computational experiments. A four-state, spatially
explicit stochastic simulation drawing its populations, geographic centroids, and commuting
flows from remote sources is specified in roughly twenty-five lines of code.

\subsection{Intra-population models: symbolic specification of disease mechanics}

\texttt{epymorph} represents disease mechanics using compartmental models in which the host
population at each node is partitioned into a set of compartments, and individuals
transition between compartments according to defined rate laws. We specify these models
symbolically. A modeler declares the compartments, declares the data and parameters the
model requires, and then defines the transitions between compartments as symbolic
expressions built from the compartment state variables and the required parameters.

As an example, the built-in SIRS model is defined in full by declaring three compartments
(\texttt{S}, \texttt{I}, \texttt{R}), the three parameters it requires (the transmission
rate $\beta$, the recovery rate $\gamma$, and the immune waning rate $\xi$), and the three transitions
between compartments:

\begin{lstlisting}[language=Python]
class SIRS(CompartmentModel):
    compartments = [
        compartment("S"),
        compartment("I"),
        compartment("R"),
    ]

    requirements = [
        AttributeDef("beta", type=float, shape=Shapes.TxN),
        AttributeDef("gamma", type=float, shape=Shapes.TxN),
        AttributeDef("xi", type=float, shape=Shapes.TxN),
    ]

    def edges(self, symbols):
        [S, I, R] = symbols.all_compartments
        [β, γ, ξ] = symbols.all_requirements
        N = S + I + R
        return [
            edge(S, I, rate=β * S * I / N),
            edge(I, R, rate=γ * I),
            edge(R, S, rate=ξ * R),
        ]
\end{lstlisting}

where N is the total population at the node. The transition rates are written as ordinary
mathematical expressions in the compartment symbols, which \texttt{epymorph} parses using
the SymPy symbolic mathematics library~\citep{meurer_sympy_2017}. We adopt this symbolic approach for three
reasons. First, the model definition reads as the mathematics it represents, which increases
transparency and reduces the gap between a published model and its implementation. Second,
because the rate laws are symbolic, \texttt{epymorph} can automatically render a diagram of
the compartmental structure for user inspection and can analyze the model to determine which
parameters and data it requires. Third, the symbolic expressions are compiled (lambdified)
into fast numerical functions at simulation time, so the readability of the specification
carries no runtime penalty.

The transition system also supports competing and branching processes. A single source
compartment may have several outgoing transitions, and a transition may be specified as a
fork, in which individuals leaving a compartment are partitioned among two or more
destinations according to a set of probabilities. For example, the built-in SIRH model
routes infectious individuals into either a hospitalized compartment or a recovered
compartment using a fork whose branch probabilities are set by a hospitalization probability
parameter. Compartments may also carry tags that modify their treatment elsewhere in the
simulation, such as marking a hospitalized compartment as immobile so that those individuals
are excluded from movement. Pseudo-compartments representing exogenous birth and death allow
models to include host demography by adding transitions to or from outside the modeled
population.

\subsection{Movement models: sub-daily tempo and movement clauses}

Host movement between nodes is described by a movement model, which divides each simulated
day into one or more sub-intervals, which we call tau steps, and specifies movement during
each step through a set of movement clauses. These tau steps are the same sub-daily
increments over which the stochastic disease process advances: \texttt{epymorph} simulates
dynamics with a tau-leaping scheme in which each step advances both movement and
transmission (detailed in ``The simulation engine'' below). Each clause computes the number
of individuals requested to move from each node to every other node at a given step,
typically as a function of node populations, inter-node distances or commuting data, and
clause-specific parameters.

% Movement is treated as requested rather than guaranteed, because the number of individuals available to move may be smaller than the number requested; in that case fewer individuals move.

As an example, the built-in centroids movement model divides the day into two tau steps of
lengths 1/3 and 2/3. During the first step a fixed proportion of each node's population is
drawn to commute, and the destinations of those commuters are assigned by a multinomial draw
against a distance-based dispersal kernel. During the second step the commuters return home.
Movement clauses specify the tau step at which movers leave and the delay after which they
return, so that models can represent movement on a range of temporal scales, from daily
commuting to longer trips. Because movement is defined independently of disease mechanics,
the same disease process can be simulated under different movement assumptions without
modification, and different subpopulations within a single experiment can follow different
movement models.

% (4) k(i, j) ∝ exp(−d(i, j)/φ),
% where d(i, j) is the great-circle distance between the centroids of nodes i and j and φ is a parameter that tunes the typical distance traveled.
% The kernel is row-normalized so that each row defines a probability distribution over destinations.

\subsection{Geographic scope and automated data acquisition}

A central practical barrier to spatial modeling is the effort required to assemble
population, movement, and surveillance data for a chosen geographic extent and resolution.
\texttt{epymorph} addresses this through two coupled abstractions. The first is the
geographic scope, which defines the set of nodes in the model. For the United States, scopes
are constructed automatically for the user at standard census hierarchies, including state,
county, census tract, and census block group, and can be specified by listing the desired
areas (for example, a set of states) or by selecting all areas at a given level. The scope
determines both the spatial resolution of the simulation and the identifiers used to align
external data.

The second abstraction is the ADRIO (Abstract Data Resource Interface Object), a component
that fetches a specific variable from an external data source for a given scope and time
frame. \texttt{epymorph} provides a library of ADRIOs for several major data sources used in
U.S. spatial epidemiology, including population and demographic tables from the American
Community Survey, commuting flows from the Census commuting flows product, employment-based
movement from the LEHD Origin--Destination Employment Statistics (LODES), climate variables
from PRISM, geographic boundaries and centroids from the Census TIGER/Line files, and
infectious-disease surveillance records distributed through federal data portals, including
age-stratified COVID-19 and influenza hospitalization time series at the state and facility
levels~\citep{harris_array_2020,bureau_american_nodate,us_census_bureau_center_for_economic_studies_us_nodate,noauthor_prism_nodate}. Additional ADRIOs load data from user-supplied CSV files or in-memory
arrays, and new remote sources are added by subclassing a documented base class, so that
data sources without a built-in interface can be incorporated with minimal code. The data
layer is therefore extensible by design rather than by request: although the built-in
library is oriented toward U.S. sources, studies in other settings use the same modeling,
inference, and forecasting machinery and supply their data through these same interfaces.

When an ADRIO runs, it retrieves the requested data at the resolution of the scope, aligns
it to the model's nodes and time frame, and caches the result so that repeated experiments
do not repeat downloads. Retrieved values are also checked and repaired against declared
expectations through a small vocabulary of fill and fix strategies, so that missing or
invalid entries are handled explicitly rather than propagating silently into a simulation.
Because an ADRIO is defined by the variable it returns and the shape that variable takes,
the same component can serve either of the two roles a spatial study requires: it may supply
a parameter of the model, or it may supply the observational data that a calibration routine
targets (see ``Parameter estimation and forecasting'' below). The data layer and the
inference layer are thus reached through one interface, and a modeler need only name the
variable and the scope.

ADRIOs are one instance of a more general abstraction that unifies the data layer of
\texttt{epymorph}. ADRIOs, parameter functions, and initializers are all implemented as
simulation functions, which are lazily evaluated against a shared context consisting of the
scope, the time frame, and other already-resolved data. A simulation function may defer to
another, requesting a value it needs, so that data dependencies are resolved automatically
as a tree. This design allows a parameter to be specified as a constant, as a function of
time or space, or as the output of an ADRIO, interchangeably and without special handling by
the rest of the system. Every data value carries an explicit shape drawn from a small set of
canonical shapes (a single scalar; a value per node, N; a value per time step, T; a
node-by-node matrix, N$\times$N; or a time-by-node matrix, T$\times$N, among others). Shapes are validated
before a simulation runs, so that dimensional mismatches between a model's requirements and
the supplied data are caught immediately rather than producing silent errors during
execution.

Initialization of the compartment populations is handled by the same simulation-function
machinery, through a family of initializers that place the population into compartments at
the start of a run. These range from fully explicit initial conditions to schemes that seed
infection at chosen or randomly selected nodes, so that stochastic ensembles can be started
from randomized initial locations and seed sizes when required.

\subsection{The simulation engine}

Given a fully specified RUME, the basic simulator advances the model forward in discrete
time and produces time series of compartment populations and transition events for every
node. The default disease dynamics are stochastic. Within each tau step, \texttt{epymorph}
evaluates the rate of each transition from the current compartment values and parameters,
and then draws the number of transition events from a Poisson distribution with mean equal
to the rate multiplied by the tau step length. For a fork, the total number of individuals
leaving the source compartment is drawn from a Poisson distribution and then partitioned
among the fork's destinations with a multinomial draw against the branch probabilities. When
several transitions draw from the same source compartment and the requested totals exceed
the number of individuals available, the competing draws are reconciled with a
hypergeometric draw so that the number of individuals leaving the compartment never exceeds
the number present. This scheme is a form of tau-leaping for stochastic compartmental
dynamics, and it preserves integer-valued, non-negative compartment populations~\citep{gillespie_exact_1977,gillespie_approximate_2001}.

The movement and disease phases alternate within the simulation clock. At each tau step the
movement model relocates individuals among nodes, those individuals are subject to the local
disease process at their current location, and individuals return to their home node
according to the delays specified by the movement clauses. Because the multiple movement
models in a multi-strata experiment may define different tau step schemes, \texttt{epymorph}
computes a single combined tau step scheme that is compatible with all of them while using
the fewest possible steps, and re-indexes each movement model onto the combined clock.
Simulation output is returned as a structured object that supports selection of
compartments, events, nodes, and time ranges, along with built-in routines for plotting time
series, drawing choropleth maps, and summarizing results as quantile tables across
stochastic replicates.

\subsection{Multi-strata model composition}

Many epidemiological questions require partitioning a population into strata, such as age
groups or risk classes or multiple species (e.g., mosquitoes and humans), that differ in
their disease parameters, their movement behavior, or both. \texttt{epymorph} builds
multi-strata models by combining IPMs together with meta-relationships. A modeler defines
each stratum as a GPM, pairing an IPM, an MM, and an initializer, and then supplies the
cross-strata transmission terms that couple the strata together. The package combines the
per-stratum compartment models into a single model whose compartments are labeled by
stratum, reconciles the strata's tau step schemes, and assembles the full set of
transitions. Cross-strata transmission is specified through meta-edges that reference
compartments from different strata, so that, for example, the force of infection on
susceptible individuals in stratum `a' can depend on the infectious individuals in every
stratum `b' through a matrix of contact parameters $\beta_{ab}$:

\begin{lstlisting}[language=Python]
def meta_edges(symbols):
    [S_a, I_a, R_a] = symbols.strata_compartments("a")
    [S_b, I_b, R_b] = symbols.strata_compartments("b")
    [beta_ab] = symbols.all_meta_requirements
    N_b = S_b + I_b + R_b
    return [
        edge(S_a, I_a, rate=beta_ab * S_a * I_b / N_b),
    ]
\end{lstlisting}

Here \texttt{S\_a} is the number of susceptible individuals in stratum a, \texttt{I\_b} and
\texttt{N\_b} are the infectious and total populations of stratum b, and \texttt{beta\_ab}
is the transmission rate from stratum b to stratum a; in the general case the force of
infection on stratum a sums one such meta-edge over every stratum b. Parameters in a
multi-strata model are addressed through a hierarchical naming scheme that identifies the
stratum, the module (IPM, MM, or initializer), and the parameter, and that supports
wildcards so that a value can be assigned to one stratum, to all strata at once, or to the
cross-strata meta-component. This composition machinery makes it straightforward to assign,
for example, an SIRS process to younger strata and an SIRH process with hospitalization to
an older stratum, or to let different strata follow different movement models, while reusing
the same component definitions.

\subsection{Parameter estimation and forecasting}

\texttt{epymorph} treats calibration and forecasting as operations on a RUME, so that the
same model that is simulated forward can be fit to data and used to generate probabilistic
forecasts without re-specification. These operations are organized as a single pipeline
architecture: a shared configuration binds a RUME to a number of stochastic realizations and
to a set of unknown (i.e., estimated) parameters, and every method in the pipeline ingests
that configuration and produces output in a common structure. A forecast simulator
propagates an ensemble of realizations forward over the model's time frame to generate
probabilistic predictions, while the calibration methods additionally assimilate observed
data as they advance. Because these methods share the same configuration and output, a
modeler can move between simulating, fitting, and forecasting a model without re-specifying
it, and can compare inference algorithms on an identical model and data set, evaluating the
effect of the algorithm itself.

Calibration is framed as a sequential state-space problem. The RUME defines the latent
process; a set of observations links surveillance data to the model through a model link
that specifies how a realization's output is sub-selected, grouped, and aggregated along the
geographic, temporal, and quantity axes to be compared against the data; and a likelihood
defines the observation model relating predicted quantities to observed counts.
\texttt{epymorph} provides Poisson, negative binomial, and Gaussian likelihoods, so that the
observation model can reflect the noise structure of the surveillance stream being
assimilated, and users may supply their own. Each parameter designated as `unknown' is given
a prior distribution and a model for how it evolves over time, for example static, Brownian,
or mean-reverting (Ornstein--Uhlenbeck) dynamics, together with an optional transform that
maps an unconstrained parameter onto a valid range, such as constraining a transmission rate
to remain positive. Unknown parameters may vary across both time and nodes. The available
parameter-dynamics families, their transforms, and the resulting process used for
strictly-positive, mean-reverting transmission-rate estimation (a Black--Karasinski process)
are detailed in Supplementary Methods~\citep{drennan_estimating_2026}.

Two families of sequential algorithms are currently implemented behind this common
interface. The first is a particle filter, which represents the posterior distribution over
states and parameters by a weighted ensemble of particles that are propagated through the
process model, weighted by their likelihood against each new observation, and resampled. The
implementation normalizes particle weights in the log domain for numerical stability and
resamples systematically~\citep{gordon_novel_1993,arulampalam_tutorial_2002}. The second algorithm is an ensemble Kalman filter, which
represents the posterior by an ensemble of realizations and updates them toward each
observation using a Kalman gain computed from the ensemble covariance~\citep{evensen_ensemble_2003}. Both filters
advance the model observation by observation, applying the update step at each point where
data are available, and both return not only the fitted trajectories but also the estimated
time- and location-varying parameter values across the ensemble. To remain effective in
high-dimensional spatial settings, where a global filter degenerates as the number of nodes
grows, both filters assimilate observations with spatial localization, so that information
is incorporated node by node~\citep{ionides_iterated_2024}.

We offer these two filtering approaches so that users can choose the one best suited to
their study, trading off the advantages and disadvantages of each. Both filters accommodate
arbitrary observation likelihoods, including the Poisson and negative binomial models
appropriate for count-based surveillance: the particle filter weights particles directly by
their likelihood, whereas the ensemble Kalman filter draws perturbed observations from the
specified likelihood and assimilates them through its update step. The filters differ in how
that update behaves. The particle filter targets the exact posterior but can degenerate in
large spatial systems, where a finite ensemble loses track of the state; the ensemble Kalman
filter instead applies a moment-based linear update computed from the ensemble covariance,
which is only an approximation of the true posterior when the observation model is
non-Gaussian but remains robust as the number of nodes grows because it can adjust
individual ensemble members toward the data.

Output from any pipeline method can be summarized across realizations and across estimated
parameters, for example as ensemble means, standard deviations, or quantiles, using a common
selection and aggregation interface. Because the particle filter and the ensemble Kalman
filter expose the same interface and operate on the same RUME, a researcher can compare
inference methods on an identical model and data set, and can carry a calibrated ensemble
directly into the forecast simulator to produce predictions with real-time updating as new
data arrive.

\subsection{Implementation and availability}

\texttt{epymorph} is implemented in Python (versions 3.11 through 3.13) and builds on the
standard scientific Python stack: NumPy for numerical arrays~\citep{harris_array_2020}, SymPy for symbolic
model specification~\citep{meurer_sympy_2017}, pandas for tabular data handling~\citep{mckinney_data_2010}, SciPy for
statistical distributions, and Matplotlib for visualization~\citep{hunter_matplotlib_2007}. Its automated data
layer additionally draws on the geospatial Python ecosystem, including geopandas and
rasterio for spatial and raster data and a Census API client for demographic and commuting
data. Numerical performance is achieved by compiling the symbolic rate laws into vectorized
numerical functions and by operating on integer-typed population arrays. The package is
released under the GNU General Public License v3 (GPL-3.0) and distributed through the
Python Package Index; each release is archived at a citable DOI, and the package is
accompanied by documentation and a collection of vignettes that demonstrate single-node
simulation, multi-node and multi-strata modeling, parameter estimation, and forecasting.
[Repository and archived-version DOI to be added.]
